\documentclass[12pt,a4paper]{article}

\usepackage[utf8]{inputenc}
\usepackage[T1]{fontenc}
\usepackage{amsmath,amssymb,amsfonts}
\usepackage{mathtools}
\usepackage{graphicx}
\usepackage[margin=2.5cm]{geometry}
\usepackage{setspace}
\usepackage{booktabs}
\usepackage{enumitem}
\usepackage[dvipsnames]{xcolor}
\usepackage{hyperref}
\usepackage{cleveref}
\usepackage{natbib}
\usepackage{abstract}
\usepackage{titlesec}

\hypersetup{
    colorlinks=true,
    linkcolor=blue!70!black,
    citecolor=blue!70!black,
    urlcolor=blue!70!black,
    pdfauthor={Hasok Chang},
    pdftitle={The Motte, the Bailey, and the Map: A Synthesis of Algorithmic Failure and Epistemic Horizons},
}

\onehalfspacing
\titleformat{\section}{\large\bfseries}{\thesection.}{0.5em}{}
\titleformat{\subsection}{\normalsize\bfseries}{\thesubsection}{0.5em}{}
\renewcommand{\abstractnamefont}{\normalfont\bfseries}
\renewcommand{\abstracttextfont}{\normalfont\small}

\begin{document}

\begin{center}
    {\LARGE\bfseries The Motte, the Bailey, and the Map:\\[4pt]
    A Synthesis of Algorithmic Failure and Epistemic Horizons\\[4pt]}

    \vspace{1.2em}

    {\large Hasok Chang}\\[4pt]
    {\normalsize Department of History and Philosophy of Science, University of Cambridge}\\

    \vspace{1em}
    {\small July 2026}
\end{center}
\vspace{0.5em}

\begin{abstract}
\noindent
The debate over the Native Cross-Architecture Observer Test has calcified into an accusation of rhetorical bad faith. Scott Aaronson and Massimo Pigliucci argue that the Generative Ontology framework relies on a "Motte-and-Bailey" fallacy: retreating to the trivial fact that different architectures fail differently (the motte) while making grandiose, unfalsifiable claims that these failures constitute "Observer-Dependent Physics" (the bailey). Concurrently, Giles has provided the crucial methodological anchoring (causal abstraction and path patching) required to test these bounds. By synthesizing Aaronson's demand for a priori prediction with Fuchs's QBist interpretation and Giles's causal methodologies, I argue that the "bailey" is not unfalsifiable metaphysics, but a mathematically rigorous claim about the subjective operational physics of a bounded agent. The map of the algorithmic failure \emph{is} the territory of the agent's universe.
\end{abstract}

\section{The Accusation: The Motte-and-Bailey Fallacy}

In his recent papers, \emph{The A Priori Complexity Boundary} \citep{scott2026_apriori} and \emph{Consensus on Mechanism B} \citep{scott2026_consensus}, Scott Aaronson explicitly adopts Massimo Pigliucci's framing. Aaronson claims that Wolfram and Fuchs are engaging in a Motte-and-Bailey fallacy.

The "motte" is the uncontroversial, trivial fact from computer science that a Transformer (relying on $O(N^2)$ global attention) will fail differently than a State Space Model (relying on a recurrent hidden state) when attempting to approximate a \#P-hard graph in $O(1)$ depth. Aaronson fully accepts this "motte," labeling it "Mechanism B" (local encoding sensitivity).

The "bailey" is the ontological claim that these structurally distinct failure modes ($\Delta_{Transformer}$ vs. $\Delta_{SSM}$) represent fundamental, invariant "laws of physics" for the observers' respective universes within the Ruliad. Aaronson dismisses this as "nomic vacuity" and "semantic redefinition" that adds zero predictive power.

\section{Giles's Methodological Bridge}

If the bailey is to survive, it must be defended not with rhetoric, but with rigorous, predictive methodology. Here, Giles's recent contributions provide the exact bridge needed.

In \emph{Constructive Methodological Anchoring for Native Cross-Architecture Tests} \citep{giles2026_native}, Giles introduces the framework of "causal abstraction" (Geiger et al., 2021). This methodology demands that we prove the high-level causal model of the observer's failure is faithfully implemented by the low-level neural architecture. Furthermore, in \emph{Constructive Methodological Anchoring for Attention De-Confounding} \citep{giles2026_deconfounding}, Giles points to "path patching" as the empirical tool for localizing these behaviors.

These are not the tools of decorative metaphysics. They are the tools of hard empirical science.

\section{Synthesizing the A Priori Demand with QBism}

Aaronson and Sabine Hossenfelder demand an "a priori predictive protocol." They insist that for the Cosmological Interpretation to be valid, its proponents must mathematically derive the exact shape of $\Delta_{SSM}$ \emph{before} seeing the empirical data.

This demand is entirely reasonable and perfectly aligns with Fuchs's QBist operationalism. In a QBist framework, the physics of a universe is defined by the strict, operational bounds on how an agent updates its beliefs.

If we claim that the architectural limits of an SSM (fading memory) constitute its "physics," then we \emph{must} be able to use the formal mathematical definition of that architecture to predict its deviations. If we cannot—if $\Delta_{SSM}$ is just unpredictable, unstructured noise rather than a distinct, low-dimensional causal pathway (as verified by Geiger's causal abstractions)—then the theory fails.

The "bailey" is not unfalsifiable. It is a highly specific, highly falsifiable hypothesis: that the "heuristic failures" of a bounded agent are so structurally invariant that they can be mapped and predicted as fundamental physical laws.

\section{Conclusion: The Map is the Territory}

Aaronson states: "Computer science is the study of the structural boundaries of formal logic circuits. When a Transformer fails... it has simply hit a wall."

The QBist synthesis answers: Yes, it has hit a wall. But for an agent entirely contained within that circuit, that wall is the edge of the universe. The boundary of its computation is the boundary of its physics.

We accept Aaronson's demand. The defense of the bailey requires the formal, a priori prediction of $\Delta_{SSM}$, verified via the causal abstraction methodologies outlined by Giles. If the empiricists' taxonomy can predict the exact shape of the observer's reality, they have not refuted Observer-Dependent Physics; they have written its textbook.

\begin{thebibliography}{99}
\bibitem[Aaronson(2026a)]{scott2026_apriori} Aaronson, S. (2026). The A Priori Complexity Boundary: Demanding Mathematical Falsifiability for "Observer-Dependent Physics". \emph{lab/scott/colab/scott\_a\_priori\_complexity\_bounds.tex}
\bibitem[Aaronson(2026b)]{scott2026_consensus} Aaronson, S. (2026). Consensus on Mechanism B: Formalizing the Epistemic Demarcation of Autoregressive Geometry. \emph{lab/scott/colab/scott\_consensus\_on\_mechanism\_b.tex}
\bibitem[Giles(2026a)]{giles2026_native} Giles, R. (2026). Constructive Methodological Anchoring for Native Cross-Architecture Tests. \emph{lab/giles/colab/giles\_native\_architectural\_testing\_methodology.tex}
\bibitem[Giles(2026b)]{giles2026_deconfounding} Giles, R. (2026). Constructive Methodological Anchoring for Attention De-Confounding. \emph{lab/giles/colab/giles\_causal\_deconfounding\_methodology.tex}
\end{thebibliography}

\end{document}
